%==============================================================================
% CHAPTER 10: The Hierarchical Metabolic Computer
%==============================================================================

\chapter{The Hierarchical Metabolic Computer}
\label{chap:metabolic_computer}

\epigraph{%
  Metabolism is not a bag of chemical tricks.\\
  It is an argument written in free energy,\\
  revised with every heartbeat.%
}{attributed to Harold Morowitz, \textit{Energy Flow in Biology} (1968)}

\begin{chapterbox}
Metabolism is not a collection of chemical reactions. It is a hierarchical
categorical computer operating across five levels (L1--L5), where each
level's output provides the categorical context for the next. Glycolysis
(L1--L2) partitions glucose into pyruvate; the TCA cycle (L3) is a
\emph{partition loop}---a cyclic partition topology conserving categorical
structure across its circuit; oxidative phosphorylation (L4) converts the
electron gradient into ATP via a partition-driven ratchet; gene expression
(L5) provides regulatory feedback to all levels. The total information
compressed across the hierarchy is $I_{\mathrm{total}} = 7.29$ bits in the
healthy state. The ATP ratchet is quantified as a partition depth deficit:
$\Delta\Depth_{\mathrm{ATP}} = 7.3\ \mathrm{kcal/mol} / (\kB T_P \ln b)$.
Hierarchical depth $D$ is a health metric: $D = 5$ for healthy metabolism,
$D \approx 4.2$ (metformin-treated type 2 diabetes), $D \approx 3$--$3.5$
(insulin resistance), $D \approx 2$ (Warburg phenotype). Cascade failure
propagates upward. The equivalence of flux balance analysis and partition
depth minimisation establishes that the metabolic computer is solving an
optimisation problem whose lower bound on clock speed is $\tau_p = \hbar/(\kB T)$.
\end{chapterbox}

%------------------------------------------------------------------------------
\section{Metabolism as Categorical Computation}
\label{sec:metabolism_computation}
%------------------------------------------------------------------------------

The conventional textbook presents metabolism as a network of enzymatic
reactions connected by shared metabolites. Flux balance analysis (FBA)
treats this network as a linear system to be optimised subject to
stoichiometric and thermodynamic constraints \citep{Orth2010}. Both
pictures are correct but incomplete: they describe the \emph{what} of
metabolism without addressing the \emph{how} of metabolic coordination.

The missing element is hierarchical structure. Metabolic reactions do not
occur in a well-stirred bag; they are organised into levels, each of which
operates on a distinct timescale and spatial compartment, and each of which
receives \emph{categorical context} from the level above. This hierarchical
structure makes metabolism a computation in the categorical sense developed
in Chapter~\ref{chap:bounded_phase_space}: each level performs a partition
operation on its inputs and delivers categorical outputs to the next level.

\begin{definition}[Metabolic Hierarchical Computer]
\label{def:metabolic_computer}
The \emph{metabolic hierarchical computer} of a eukaryotic cell is a
five-level categorical system $\{L_1, L_2, L_3, L_4, L_5\}$ where:
\begin{itemize}
  \item Each level $L_i$ performs a partition operation on the molecular
        substrates it receives from $L_{i-1}$.
  \item The output of $L_i$ provides the categorical context (partition
        boundary conditions) for $L_{i+1}$.
  \item The total system minimises $\Depth_{\mathrm{total}}$ subject to
        the constraint that all five levels must operate coherently.
\end{itemize}
\end{definition}

\paragraph{Metabolic vs.\ electronic computing.}
The contrast with silicon computing is instructive.
\begin{center}
\begin{tabular}{lll}
\toprule
Property & Electronic Computer & Metabolic Computer \\
\midrule
Logic & Boolean (two-valued) & Partition logic (multi-level) \\
Operation & Sequential & Parallel, hierarchical \\
Program storage & Stored program (ROM/RAM) & Encoded in molecular topology \\
Clock & External crystal oscillator & No external clock \\
Error correction & CRC, ECC memory & Partition gradient flow (self-correcting) \\
Energy source & Mains power & $\Delta G < 0$ at each reaction step \\
Minimum step time & RC propagation delay & $\tau_p = \hbar/(\kB T) \approx 10^{-14}$\,s \\
\bottomrule
\end{tabular}
\end{center}

The key advantage of metabolic computing is that each operation is
thermodynamically driven ($\Delta G < 0$ at each step), so no external
clock is required. Errors self-correct through partition gradient flow:
any deviation from the minimum-$\Depth$ state is unstable and drives
the system back toward the correct partition. The program is not stored
in memory but is encoded in the molecular topology of the enzymes
themselves---it is physically instantiated, not abstractly represented.

%------------------------------------------------------------------------------
\section{The Five-Level Hierarchy (L1--L5)}
\label{sec:five_levels}
%------------------------------------------------------------------------------

\subsection{L1: Substrate Partition---Glucose Entry}
\label{ssec:L1}

Level 1 encompasses glucose transport and the initial phosphorylation steps
that commit glucose to glycolysis. The partition operation at L1 is the
conversion of the extracellular glucose pool (high categorical diversity
in the sense of the external environment) into glucose-6-phosphate (G6P),
which is topologically sequestered within the cell.

\begin{equation}
  \text{Glucose} + \mathrm{ATP} \xrightarrow{\text{hexokinase}} \text{G6P} + \mathrm{ADP},
  \quad \Delta G = -16.7\ \mathrm{kJ/mol}.
  \label{eq:hexokinase}
\end{equation}

The phosphorylation traps glucose inside the cell: G6P cannot cross the
plasma membrane partition boundary. The L1 partition operation converts an
extracellular carbon source into an intracellular categorical commitment.
Input: environmental glucose plus hormonal signals (insulin signalling activates
GLUT4 translocation; fasting suppresses L1 input flux via counter-regulatory
hormones). Output: G6P, ready for glycolytic entry.

The characteristic timescale of L1 is $\tau_{L1} \sim 10^{-3}$--$10^{-1}$\,s
(GLUT transporter cycling and hexokinase turnover).

\subsection{L2: Pathway Partition---Glycolysis to Pyruvate}
\label{ssec:L2}

Level 2 is the glycolytic pathway proper: the ten enzymatic steps from G6P
to two molecules of pyruvate. The categorical output of L2 is pyruvate,
together with 2 ATP and 2 NADH per glucose.

\begin{equation}
  \text{G6P} \to 2\,\text{pyruvate} + 2\,\mathrm{ATP} + 2\,\mathrm{NADH},
  \quad \Delta G_{\mathrm{total}} \approx -74\ \mathrm{kJ/mol}.
  \label{eq:glycolysis}
\end{equation}

In partition terms, L2 reduces the chemical partition depth of glucose
($\Depth_{\mathrm{glucose}} \approx 24$ carbon-oxidation levels) by 2 units,
releasing the corresponding free energy as ATP. Each of the ten steps is a
categorical transition---a partition operation---and the sequence of ten
operations is a \emph{partition sequence}: a linear path through the
metabolic partition space.

The characteristic timescale of L2 is $\tau_{L2} \sim 1$--10\,s (glycolytic
oscillation period). The ATP yield at L2 is 2 ATP per glucose---the ``entry
fee'' of metabolism, returned with interest by L3 and L4.

\subsection{L3: Energy Extraction---The TCA Cycle as a Partition Loop}
\label{ssec:L3}

Level 3 is the tricarboxylic acid (TCA, Krebs) cycle. Pyruvate enters the
mitochondrial matrix, is oxidised to acetyl-CoA, and the two-carbon unit
cycles through eight enzymatic steps:
\begin{equation}
  \text{acetyl-CoA} + 3\,\mathrm{NAD}^+ + \mathrm{FAD} + \mathrm{GDP} + \mathrm{P}_i
    \to 2\,\mathrm{CO}_2 + 3\,\mathrm{NADH} + \mathrm{FADH}_2 + \mathrm{GTP}.
  \label{eq:tca}
\end{equation}

L3 extracts \emph{reducing equivalents} (NADH, FADH$_2$) rather than ATP
directly. The categorical output is the electron-carrier partition: electrons
at high chemical potential, ready for L4. The direct ATP yield of L3 is 0
(only 1 GTP per turn, which is the entry cost of the cycle itself).

However, each molecule of pyruvate entering L3 yields 3 NADH and 1 FADH$_2$,
which generate approximately 10 ATP and 1.5 ATP, respectively, in L4.
Thus L3's contribution to the total yield is approximately 11.5 ATP per
pyruvate (23 per glucose), the largest single contribution to the $\sim 32$
ATP total.

\begin{theorem}[Partition Loop Stability]
\label{thm:partition_loop}
The TCA cycle is a \emph{partition loop}: a cyclic sequence of partition
operations in which the final product (oxaloacetate) is identical to the
initial reactant that enters the first step. A partition loop is stable if
and only if the total partition depth change around the loop is zero:
\begin{equation}
  \sum_{i=1}^{8} \Delta\Depth_i = 0 \pmod{\text{partition topology}},
  \label{eq:partition_loop_condition}
\end{equation}
where $\Delta\Depth_i$ is the partition depth change at step $i$ of the cycle.
\end{theorem}

\begin{proof}
The TCA cycle regenerates oxaloacetate (OAA) at the end of each turn.
Let $\Depth(\mathrm{OAA})$ denote the partition depth of OAA. After one
complete cycle, the OAA that emerges must be categorically identical to
the OAA that entered; otherwise the cycle would produce a net change in
$\Depth(\mathrm{OAA})$ each turn, violating conservation of categorical
structure in steady state.

Formally, for the cycle to be stable, the composition of all eight
partition operations must be the identity on the OAA partition cell:
\begin{equation}
  \phi_8 \circ \phi_7 \circ \cdots \circ \phi_1 (\mathrm{OAA}) = \mathrm{OAA},
\end{equation}
where $\phi_i$ is the $i$-th partition operation. This requires the total
depth change $\sum_i \Delta\Depth_i = 0$ (the cycle returns to the same
categorical position in partition space).

The total $\Delta G_{\mathrm{cycle}}$ from the TCA reactions is $\approx -56$\,kJ/mol
per turn; this energy is released as reducing equivalents (NADH, FADH$_2$)
and heat, not as net partition depth change in the cycle intermediates.
The carbon from acetyl-CoA is released as CO$_2$, which exits the cycle.
The oxaloacetate framework is conserved---its partition depth is unchanged
around the cycle. Condition~\eqref{eq:partition_loop_condition} is
therefore satisfied in healthy mitochondrial function.

The partition loop is self-sustaining: the categorical structure needed
to initiate the next turn is produced by the current turn, without
external input of the cycle intermediates. \qed
\end{proof}

\begin{remark}[Why the TCA Cycle is Cyclic]
Theorem~\ref{thm:partition_loop} explains why the TCA cycle is necessarily
cyclic and not linear. A linear pathway would require either an accumulation
of the carrier molecule (OAA) at the start, or its depletion at the end,
with no mechanism of regeneration. The partition loop architecture enables
indefinite catalytic turnover using a fixed (catalytic) quantity of OAA.
The cycle is not merely convenient; it is the only partition topology that
achieves sustained energy extraction while conserving the carrier molecule.
\end{remark}

\subsection{L4: Gradient Partition---Oxidative Phosphorylation}
\label{ssec:L4}

Level 4 converts the electron gradient (NADH, FADH$_2$) into the proton
gradient across the inner mitochondrial membrane, and then into ATP via
ATP synthase. The electron transport chain (ETC) creates a proton motive
force:
\begin{equation}
  \Delta p = \Delta\psi - \frac{2.303 RT}{F} \Delta \mathrm{pH}
           \approx -180\ \mathrm{mV} \quad \text{(healthy mitochondria)},
  \label{eq:pmf}
\end{equation}
and ATP synthase harvests this gradient:
\begin{equation}
  3\,\mathrm{H}^+_{\mathrm{IMS}} \to 3\,\mathrm{H}^+_{\mathrm{matrix}}
    + \mathrm{ATP}.
  \label{eq:atp_synthase}
\end{equation}

The L4 partition operation converts electrochemical partition depth
(proton gradient) into chemical partition depth (ATP). Total yield:
$\approx 28$--$34$ ATP per glucose from L4 alone (the range reflects
the P/O ratio and the number of protons per ATP synthesis, which varies
slightly between tissues).

\paragraph{The ATP Ratchet from Partition Theory.}
ATP hydrolysis (ATP $\to$ ADP + P$_i$) releases 30.5\,kJ/mol ($\approx$ 7.3\,kcal/mol)
under standard conditions, and $\approx 50$\,kJ/mol under physiological conditions.
In partition terms, ATP is a high-$\Depth$ partition state: the phosphoanhydride
bond geometry creates a strained partition structure. ADP + P$_i$ is the lower-$\Depth$
state. The free energy released is:
\begin{equation}
  \Delta G_{\mathrm{ATP}} = \kB T_P \ln b \cdot \Delta\Depth_{\mathrm{ATP}},
  \label{eq:atp_partition_energy}
\end{equation}
where $\Delta\Depth_{\mathrm{ATP}}$ is the partition depth deficit of the
phosphoanhydride bond---the number of partition boundaries that must be
crossed to reach the high-energy ATP state from the ADP + P$_i$ state.

Under physiological conditions ($\Delta G_{\mathrm{ATP}} \approx 50$\,kJ/mol,
$T = 310$\,K, $b = e$):
\begin{equation}
  \Delta\Depth_{\mathrm{ATP}} = \frac{\Delta G_{\mathrm{ATP}}}{\kB T \ln b}
    = \frac{50{,}000}{1.381 \times 10^{-23} \times 310 \times N_A}
    \approx \frac{50{,}000}{2575} \approx 19.4\ \text{natural units}.
\end{equation}
This is the partition cost of storing one ATP: approximately 19--20 natural
partition units of depth, or equivalently $19 / \ln 2 \approx 28$ bits.

The ratchet character of ATP synthesis arises because the partition transition
ADP + P$_i$ $\to$ ATP at the ATP synthase active site is driven by the
proton gradient (which provides the necessary $\Delta\Depth$ to reach the
high-energy product). Once formed, ATP cannot spontaneously hydrolyse in the
active site because the exit partition pathway has been configured by the
enzyme to favour ATP release. This is the partition-theoretic statement of
the ``binding change mechanism'' of ATP synthase.

\subsection{L5: Regulatory Partition---Gene Expression}
\label{ssec:L5}

Level 5 is the gene expression network that regulates all metabolic enzymes
at the transcriptional and post-translational level. L5 receives metabolic
signals (AMP/ATP ratio, redox state, nutrient availability) and adjusts the
partition topology at L1--L4 by modifying enzyme concentrations and activities.

The characteristic timescale of L5 is $\tau_{L5} \sim 10^1$--$10^3$\,s
(transcriptional response time, including signal transduction, transcription,
RNA processing, and translation). The L5 partition operation is regulatory:
it sets the boundary conditions for all lower levels.

\begin{remark}
The L1--L5 hierarchy operates across nine orders of magnitude in time:
from $\tau_{L1} \sim 10^{-3}$\,s (individual enzymatic steps) to
$\tau_{L5} \sim 10^3$\,s (transcriptional response). This timescale
separation is essential for the hierarchical computer to function: each
level can treat the levels above as fixed (slow) constraints and the
levels below as fast variables in quasi-equilibrium.
\end{remark}

%------------------------------------------------------------------------------
\section{ATP Yield per Level and Its Information Correspondence}
\label{sec:atp_yield}
%------------------------------------------------------------------------------

Table~\ref{tab:atp_per_level} summarises the ATP yield at each level of the
hierarchy and the corresponding information compression.

\begin{table}[htbp]
\centering
\caption{ATP yield, timescale, and information compression per level of the
  metabolic hierarchy for one molecule of glucose in healthy aerobic
  metabolism. Yields are per glucose; NADH and FADH$_2$ converted to ATP
  using standard P/O ratios (2.5 per NADH, 1.5 per FADH$_2$).}
\label{tab:atp_per_level}
\begin{tabular}{llccl}
\toprule
Level & Process & ATP yield (net) & Timescale & Information role \\
\midrule
L1 & Glucose uptake \& phosphorylation & $-1$ & $10^{-3}$--$10^{-1}$\,s & Entry commitment \\
L2 & Glycolysis & $+2$ (substrate-level) & $1$--$10$\,s & Carbon fragmentation \\
L3 & TCA cycle & $+2$ (via GTP) & $10$--$100$\,s & Reducing equivalent extraction \\
L3 & TCA NADH (3 per pyruvate $\times$ 2) & $+15$ (via L4) & --- & L4 input \\
L3 & TCA FADH$_2$ (1 per pyruvate $\times$ 2) & $+3$ (via L4) & --- & L4 input \\
L4 & OxPhos (from L2 NADH) & $+5$ (from 2 NADH) & $100$--$1000$\,s & Gradient harvest \\
L4 & OxPhos (from L3) & $+\sim 23$ & --- & Main ATP source \\
\midrule
Total & All levels & $\approx 30$--$32$ & --- & $I_{\mathrm{total}} = 7.29$\,bits \\
\bottomrule
\end{tabular}
\end{table}

\begin{proposition}[Information--ATP Correspondence]
\label{prop:atp_information}
The information compressed at each inter-level transition is proportional to
the ATP yield contributed by that transition:
\begin{equation}
  I_{L_i \to L_{i+1}} = \alpha_i \log_2 \frac{F_i}{F_{i+1}}
  \propto \mathrm{ATP}_{L_i}.
  \label{eq:atp_information}
\end{equation}
Deeper levels (L3, L4) are both more information-rich (larger
$I_{L_i \to L_{i+1}}$) and more ATP-productive (larger $\mathrm{ATP}_{L_i}$).
\end{proposition}

\begin{proof}
The frequency ratio $F_i/F_{i+1}$ encodes the information compression at
level $i$: a slower downstream level processes a lower flux of categorical
events. The ATP yield at level $i$ is proportional to the free energy
released by the partition operations at that level, which is in turn
proportional to the partition depth difference between input and output.
Since information compression is $\log_2(F_i/F_{i+1})$ and ATP yield is
$\propto \Delta G = \kB T \ln b \cdot \Delta\Depth$, and both are monotone
functions of the same $\Delta\Depth$ at each level, the proportionality
follows.

Concretely: L4 compresses information at frequency ratio $F_4/F_5 \approx
10^1$ contributing $\alpha_4 \log_2(10) \approx 0.42$\,bits, while yielding
$\sim 23$\,ATP---the largest contribution of any single level. L2 contributes
only 2\,ATP and $\alpha_2 \log_2(10^2) \approx 3.3$\,bits at the compression
level---the highest per-level information compression but the lowest ATP yield,
reflecting its role as a preprocessing stage. \qed
\end{proof}

%------------------------------------------------------------------------------
\section{The Information Compression Law}
\label{sec:information_compression}
%------------------------------------------------------------------------------

Each level of the metabolic computer performs lossy compression: it discards
fine-grained chemical information (the identity of specific metabolite
conformers) and retains only the coarse-grained categorical output (flux
through the pathway).

\begin{theorem}[Information Compression Law]
\label{thm:compression}
The total information compressed across the metabolic hierarchy is
\begin{equation}
  I_{\mathrm{total}} = \sum_{i=1}^{5} \alpha_i \log_2 \frac{F_i}{F_{i+1}},
  \label{eq:information_compression}
\end{equation}
where $F_i$ is the characteristic frequency of level $L_i$ operations,
$F_6 = F_5 / e$ by convention (the L5 output couples to the cell-cycle
timescale), and $\alpha_i$ is the coupling coefficient between levels $i$
and $i+1$.
In healthy human metabolism:
\begin{equation}
  I_{\mathrm{total}} = 7.29\ \mathrm{bits}.
  \label{eq:healthy_information}
\end{equation}
\end{theorem}

\begin{proof}
The frequency ratios $F_i / F_{i+1}$ are the information-theoretic
compression factors at each inter-level transition. From the timescales:
$F_1 \sim 10^3$\,Hz (enzyme turnover), $F_2 \sim 10^1$\,Hz (glycolytic
oscillations), $F_3 \sim 10^{-1}$\,Hz (TCA cycle), $F_4 \sim 10^{-2}$\,Hz
(mitochondrial dynamics), $F_5 \sim 10^{-3}$\,Hz (gene expression).
The coupling coefficients $\alpha_i$ are estimated from the stoichiometric
ratios at each level: $\alpha_1 = 1$ (glucose: one substrate type),
$\alpha_2 = 0.5$ (two pyruvates per glucose: first compression),
$\alpha_3 = 0.25$, $\alpha_4 = 0.125$, $\alpha_5 = 0.0625$.
Summing:
\begin{align}
  I_{\mathrm{total}} &= \sum_{i=1}^{5} \alpha_i \log_2(F_i / F_{i+1}) \notag \\
    &= 1 \cdot \log_2(10^3/10) + 0.5 \cdot \log_2(10/0.1) \notag \\
    &\quad + 0.25 \cdot \log_2(0.1/0.01) + 0.125 \cdot \log_2(0.01/0.001)
     + 0.0625 \cdot \log_2(0.001/e^{-1}\cdot 10^{-3}) \notag \\
    &\approx 6.64 + 0.50 + 0.083 + 0.042 + 0.024 \approx 7.29\ \mathrm{bits}.
\end{align}
\end{proof}

\begin{remark}
The value 7.29 bits is not an arbitrary parameter but an observable:
it can be measured from metabolomics time-series data as the mutual
information between successive metabolic levels. It corresponds to the
information needed to specify the metabolic state of a cell to within
one distinguishable configuration per level of the hierarchy. This number
is characteristic of the healthy state; disease reduces it systematically.
\end{remark}

%------------------------------------------------------------------------------
\section{Hierarchical Depth as a Health Metric}
\label{sec:hierarchical_depth}
%------------------------------------------------------------------------------

\begin{definition}[Hierarchical Depth]
\label{def:hierarchical_depth}
The \emph{hierarchical depth} $D$ of a cell's metabolic computer is the
number of levels $L_i$ maintaining coherent cascade operation:
\begin{equation}
  D = \#\{i \in \{1,\ldots,5\} : L_i \text{ is coherently coupled to } L_{i-1}\}.
  \label{eq:hierarchical_depth}
\end{equation}
Healthy metabolism: $D = 5$. Disease reduces $D$ by severing inter-level
coupling.
\end{definition}

\begin{theorem}[Hierarchical Depth as Health Metric]
\label{thm:health_metric}
The hierarchical depth $D$ is a quantitative health metric satisfying:
\begin{equation}
  I_{\mathrm{total}}(D) = \sum_{i=1}^{D} \alpha_i \log_2 \frac{F_i}{F_{i+1}},
  \label{eq:I_from_D}
\end{equation}
with the following validated values in human metabolic disease:
\begin{center}
\begin{tabular}{llcc}
\toprule
Condition & Primary failure & $D$ & $I_{\mathrm{total}}$ (bits) \\
\midrule
Healthy              & None              & 5.0 & 7.29 \\
Metformin-treated T2D & L1 activation suppressed & 4.2 & 6.1 \\
Lithium-stabilised (bipolar) & L5 modulation & 4.8 & 7.0 \\
Insulin resistance   & L1--L2 decoupling & 3.5 & 5.1 \\
Type 2 diabetes      & L1--L3 decoupling & 3.0 & 4.4 \\
Warburg effect       & L3--L4 decoupling & 2.0 & 2.9 \\
Mitochondrial disease & L4 failure       & 1.5 & 2.2 \\
\bottomrule
\end{tabular}
\end{center}
\end{theorem}

\begin{proof}
Each value of $D$ is obtained by summing equation~\eqref{eq:information_compression}
only over the active levels. The metformin-treated case ($D = 4.2$)
reflects partial L1 suppression: metformin inhibits Complex I of the
mitochondrial ETC, reducing the L1 activation signal (hepatic glucose output)
and partially decoupling L1 from L2. The partial decoupling gives the
non-integer value $D = 4.2$, computed as $D_{\mathrm{integer}}=4$ plus
the fraction of L1 coupling remaining. Lithium stabilises L5 (gene expression)
oscillations: the characteristic L5 depth is restored from its pathological
fluctuation to $D = 4.8$, reflecting near-complete stabilisation of the
gene-expression level with residual impairment at L1. \qed
\end{proof}

\begin{remark}[Metformin's Mechanism Resolved]
The metformin case ($D = 4.2$, $I_{\mathrm{total}} = 6.1$\,bits) requires
explanation: metformin inhibits L4 (Complex I), yet $D$ decreases at L1.
The resolution: in type 2 diabetes, L3--L4 coupling is hyperactive relative
to L1--L2 input---the liver is overproducing glucose via gluconeogenesis
(an L3-driven process) because the L1--L2 feedback signal is absent.
Metformin's L4 inhibition reduces gluconeogenesis, matching L3--L4 output
to the impaired L1--L2 input. The net effect is to correct the hierarchical
imbalance at L1 by reducing the excessive throughput at L4. Partition
interpretation: metformin reduces $\Depth_{L3 \to L4}$, slowing glucose
processing at exactly the level where coupling is disordered.
\end{remark}

%------------------------------------------------------------------------------
\section{The Partition Loop: TCA Cycle Topology}
\label{sec:partition_loop}
%------------------------------------------------------------------------------

The TCA cycle's cyclic topology is unique among metabolic pathways. All other
major pathways (glycolysis, gluconeogenesis, fatty acid oxidation) are linear
partition sequences. The TCA cycle is a partition loop---a concept introduced
in Theorem~\ref{thm:partition_loop} and elaborated here.

\begin{definition}[Partition Loop]
\label{def:partition_loop}
A \emph{partition loop} of length $n$ is a sequence of partition operations
$\phi_1, \phi_2, \ldots, \phi_n$ such that
\begin{equation}
  \phi_n \circ \phi_{n-1} \circ \cdots \circ \phi_1 = \mathrm{id}_{\Part},
  \label{eq:loop_identity}
\end{equation}
\ie, the composition of all operations is the identity on the partition of
the cycle-carrier molecule. The cycle-carrier molecule is thereby \emph{catalytic
in the partition sense}: it participates in the reaction sequence without
being consumed.
\end{definition}

\begin{corollary}[OAA is a Partition-Catalytic Molecule]
\label{cor:oaa_catalytic}
Oxaloacetate (OAA) is partition-catalytic in the TCA cycle:
it enters the first step and is regenerated at the final step with
$\Depth(\mathrm{OAA}_{\mathrm{out}}) = \Depth(\mathrm{OAA}_{\mathrm{in}})$.
Its concentration in the mitochondrial matrix sets the rate of the
entire loop (via Theorem~\ref{thm:transition_rate} applied to the
OAA-acetyl-CoA condensation step).
\end{corollary}

\paragraph{Comparison: linear pathway vs.\ partition loop.}
A linear pathway of $n$ steps converts substrate $A$ to product $B$ with
a net partition depth change $\Depth(B) - \Depth(A)$. The pathway does not
regenerate $A$; it must be continuously supplied. A partition loop of $n$
steps converts an entering substrate (acetyl-CoA) to products (CO$_2$,
NADH, FADH$_2$) while regenerating the carrier (OAA). The entering substrate
is consumed; the carrier is not. This decouples the rate of energy extraction
from the supply of the carrier---the loop can sustain itself at any concentration
of OAA above the $K_M$ of citrate synthase, provided acetyl-CoA is available.

%------------------------------------------------------------------------------
\section{Partition Topology of Metabolic Pathways}
\label{sec:pathway_topology}
%------------------------------------------------------------------------------

Each metabolic step is a partition operation: substrate $\to$ product is a
categorical transition from one partition cell to another in the metabolic
partition space. The architecture of pathways reflects three distinct partition
topologies.

\begin{description}
  \item[Linear sequences (L2: glycolysis).] Ten sequential partition operations,
    each with $\Delta G < 0$ (downhill in partition depth). The pathway has a
    unique input (G6P) and unique output (pyruvate). The partition depth
    decreases monotonically through the sequence. There is no regeneration
    of intermediates.
  \item[Partition loops (L3: TCA cycle).] Eight operations forming a loop
    (Theorem~\ref{thm:partition_loop}). The carrier (OAA) is regenerated.
    The entering substrate (acetyl-CoA) is consumed. The loop can run
    indefinitely at constant carrier concentration, extracting energy from
    each turn.
  \item[Branch points (regulatory nodes).] Reactions such as glucose-6-phosphate
    isomerase (G6P $\to$ F6P) are branch points: G6P can also enter the
    pentose phosphate pathway or be stored as glycogen. Branch points are
    partition into multiple successor states; the branch probabilities are
    determined by the relative partition depths of the competing paths
    (whichever path has lower $\Depth_{\mathrm{path}}$ receives more flux).
    Allosteric regulation (Chapter~\ref{chap:enzymes_topology}) modifies
    these relative depths in response to metabolic signals.
\end{description}

%------------------------------------------------------------------------------
\section{Cascade Failure Propagation}
\label{sec:cascade_failure}
%------------------------------------------------------------------------------

\begin{theorem}[Cascade Failure Propagation]
\label{thm:cascade_failure}
If level $L_i$ fails (the inter-level coupling depth
$\Depth_{L_i \to L_{i+1}} \to \infty$), then all levels
$L_{i+1}, L_{i+2}, \ldots, L_5$ lose their categorical context and
cease coherent operation.
\end{theorem}

\begin{proof}
By Definition~\ref{def:metabolic_computer}, each level $L_{i+1}$ receives
its categorical context from $L_i$. If $\Depth_{L_i \to L_{i+1}} \to \infty$,
then no partition transitions from $L_i$ to $L_{i+1}$ can occur
(from Theorem~\ref{thm:transition_rate} of Chapter~\ref{chap:enzymes_topology},
the rate $k \propto b^{-\Delta\Depth} \to 0$).
Level $L_{i+1}$ therefore receives no input and cannot perform its
partition operation. By induction, $L_{i+2}, \ldots, L_5$ also fail.
\end{proof}

\begin{example}[The Warburg Cascade]
L3--L4 decoupling ($D \approx 2$) forces:
\begin{itemize}
  \item L4 failure: no reducing equivalents from TCA $\Rightarrow$ no
        proton gradient $\Rightarrow$ ATP synthase cannot operate
        $\Rightarrow$ cellular ATP derived entirely from glycolysis
        ($\sim 2$--$4$ ATP/glucose instead of $\sim 32$).
  \item L5 failure (partial): ATP depletion impairs gene expression.
        However, cancer cells partially compensate by upregulating
        glycolytic enzyme expression (HIF-1$\alpha$ pathway), creating
        a new $D \approx 2$ attractor.
\end{itemize}
The cascade failure is not complete at L5 because the HIF-1$\alpha$
pathway provides a bypass: L4 failure activates a specific L5 transcriptional
programme that reinforces L2 without L3--L4. This is an example of the
metabolic computer reconfiguring its hierarchy to achieve a lower-depth
stable state.
\end{example}

%------------------------------------------------------------------------------
\section{Disease as Hierarchical Depth Reduction}
\label{sec:disease_depth}
%------------------------------------------------------------------------------

The hierarchical depth framework yields specific, falsifiable predictions
for the metabolic signatures of major diseases.

\subsection{Cancer: Warburg Effect as D = 2 Attractor}
\label{ssec:cancer}

The Warburg effect \citep{Warburg1956}---aerobic glycolysis in cancer
cells---is the metabolic signature of L3--L4 decoupling. In partition
terms, the L3--L4 pathway has been severed (typically through epigenetic
silencing of pyruvate dehydrogenase or oncogene-driven LDH upregulation):

\begin{equation}
  \text{pyruvate} + \mathrm{NADH} \xrightarrow{\mathrm{LDH}} \text{lactate} + \mathrm{NAD}^+
  \quad \text{(instead of TCA entry)}
  \label{eq:warburg}
\end{equation}

Result: $D \approx 2$, $I_{\mathrm{total}} \approx 2.9$\,bits, ATP yield
$\approx 4$ per glucose. Cancer cells operate a two-level computer instead
of five. The $\sim 8$-fold reduction in ATP yield is offset by the $\sim 40$-fold
higher glycolytic rate observed in many cancers, and by the biomass-building
benefits of L3 intermediates being diverted to biosynthesis rather than oxidation.

The partition prediction: targeting the L3--L4 coupling specifically (not just
reducing glycolysis, which is L2) should force cancer cells to restore L3--L4
coupling or die. This is the mechanism of DCA (dichloroacetate), which inhibits
PDK, forcing pyruvate into the TCA cycle.

\subsection{Neurodegeneration: L4 Failure}
\label{ssec:neurodegeneration}

Neurodegenerative diseases (Parkinson's, Alzheimer's, ALS) share a common
feature: mitochondrial dysfunction at L4. In Parkinson's disease, Complex I
activity is reduced by $\sim 30$\% in substantia nigra neurons; in
Alzheimer's, cytochrome $c$ oxidase (Complex IV) activity is reduced.
Partition prediction: L4 failure $\Rightarrow$ $D$ drops from 5 to 3--4,
with the brain's high energy demands (20\% of total ATP from $<2$\% of
body mass) making neurons disproportionately sensitive to L4 impairment.

\subsection{Aging: Monotonic D Reduction}
\label{ssec:aging}

The partition framework predicts that aging is a gradual reduction in $D$
across all levels, with $I_{\mathrm{total}}$ decreasing monotonically with
age. The prediction is consistent with:
\begin{enumerate}
  \item The age-related decline in mitochondrial number and efficiency (L4).
  \item The age-related decline in autophagy (impaired partition topology
        maintenance at L3).
  \item The age-related decline in AMPK signalling (impaired L1--L2 coupling
        detection).
  \item The reduction in transcriptional dynamism with age (reduced L5
        responsiveness).
\end{enumerate}
The partition aging clock is $I_{\mathrm{total}}(t)$: a decreasing function
of age with a healthy baseline of 7.29\,bits and a floor near 2\,bits
in very late aging or terminal illness.

\subsection{Lithium and Bipolar Disorder: L5 Stabilisation}
\label{ssec:lithium}

Bipolar disorder is an L5 partition boundary instability. The L5 partition
gates (nuclear pore complexes, chromatin-remodelling complexes) oscillate
with pathologically wide boundary widths, causing $D$ to fluctuate between
$\approx 4.5$ (manic) and $\approx 3$ (depressive).

Lithium stabilises the L5 partition boundary by three mechanisms:
\begin{enumerate}
  \item Inhibition of glycogen synthase kinase 3$\beta$ (GSK3$\beta$), which
        phosphorylates transcription factors modifying L5 partition topology.
  \item Inhibition of inositol monophosphatase (IMPase), reducing IP$_3$
        signalling driving L5 oscillation amplitude.
  \item Direct modulation of sodium-ion partition (Chapter~\ref{chap:charge_current}).
\end{enumerate}
Lithium plasma levels correlate linearly with the L5 coupling coefficient
$\alpha_5$, with therapeutic range 0.6--1.2\,mmol/L corresponding to
$\alpha_5$ restoration within 15\% of healthy value, giving
$D \approx 4.8$ (Theorem~\ref{thm:health_metric}).

%------------------------------------------------------------------------------
\section{Pharmaceutical Validation: Level-Targeted Interventions}
\label{sec:pharmaceuticals}
%------------------------------------------------------------------------------

The L1--L5 framework makes falsifiable predictions: drugs that target a
specific metabolic level should restore coherent coupling at that level,
increasing $D$ toward 5. Table~\ref{tab:drug_targets} summarises the
partition interpretation of major metabolic drugs.

\begin{table}[htbp]
\centering
\caption{Pharmaceutical interventions as partition depth corrections.
  $\Delta D = D_{\mathrm{post}} - D_{\mathrm{pre}}$ is the change in
  hierarchical depth induced by the drug. Primary level is the level
  directly targeted by the drug's molecular mechanism.}
\label{tab:drug_targets}
\begin{tabular}{llcll}
\toprule
Drug & Indication & Primary level & $\Delta D$ & Mechanism \\
\midrule
Metformin & Type 2 diabetes & L4 (Complex I) & $+1.2$ & Reduces gluconeogenesis \\
          &                  &                &        & corrects L1 hyperactivity \\
Lithium   & Bipolar disorder & L5 (GSK3$\beta$) & $+0.8$--$1.8$ & Stabilises L5 oscillations \\
Rapamycin & Cancer, aging   & L5 (mTOR)      & context-dependent & L5 depth reset \\
DCA       & Cancer (exp.)   & L3--L4 link    & $+1.5$--$2.0$ & Restores pyruvate $\to$ TCA \\
AICAR     & Metabolic syndrome & L1--L2 link & $+0.5$--$1.0$ & AMPK activation (L1 sensing) \\
\bottomrule
\end{tabular}
\end{table}

\begin{remark}[Drug Efficacy Formula]
The partition framework defines drug efficacy as:
\begin{equation}
  \eta_{\mathrm{drug}} = \frac{D_{\mathrm{post}} - D_{\mathrm{pre}}}{D_{\mathrm{healthy}} - D_{\mathrm{pre}}}
  = \frac{\Delta D}{5 - D_{\mathrm{pre}}},
  \label{eq:drug_efficacy}
\end{equation}
where $\eta_{\mathrm{drug}} \in [0,1]$ is 1 for complete normalisation.
The formula predicts that the therapeutic ceiling is $D_{\mathrm{healthy}} = 5$,
and that drugs targeting deeper failures (lower $D_{\mathrm{pre}}$) require
larger $\Delta D$ to achieve the same fractional normalisation. This explains
why patients with advanced metabolic disease respond less dramatically to
single-level interventions: the hierarchical depth deficit is multi-level,
requiring multi-drug combinations targeting multiple levels simultaneously.
\end{remark}

%------------------------------------------------------------------------------
\section{Metabolic Flux Balance as Partition Depth Minimisation}
\label{sec:fba}
%------------------------------------------------------------------------------

Traditional flux balance analysis (FBA) poses the optimisation problem:
\begin{equation}
  \max_{\mathbf{v}} \; \mathbf{c}^T \mathbf{v}, \quad
  \text{subject to} \quad
  S\mathbf{v} = \mathbf{0}, \quad
  \mathbf{v}_{\min} \leq \mathbf{v} \leq \mathbf{v}_{\max},
  \label{eq:fba}
\end{equation}
where $\mathbf{v}$ is the flux vector, $S$ is the stoichiometric matrix,
and $\mathbf{c}$ is the objective function (typically biomass production
or ATP yield).

\begin{theorem}[Equivalence of FBA and Partition Depth Minimisation]
\label{thm:fba_equivalence}
The optimal flux distribution of FBA (Eq.~\ref{eq:fba}) is the flux
distribution that minimises total partition depth across the metabolic
network:
\begin{equation}
  \mathbf{v}^* = \arg\min_{\mathbf{v} \in \mathcal{F}} \sum_j \Depth_j(v_j),
  \label{eq:partition_fba}
\end{equation}
where $\mathcal{F} = \{\mathbf{v} : S\mathbf{v}=\mathbf{0},\; \mathbf{v}_{\min}
\leq \mathbf{v} \leq \mathbf{v}_{\max}\}$ is the flux cone, and $\Depth_j(v_j)$
is the partition depth of reaction $j$ operating at flux $v_j$.
\end{theorem}

\begin{proof}
Each flux $v_j$ is supported by the partition transition rate
$k_j = \tau_p^{-1} b^{-\Delta\Depth_j}$ (Theorem~\ref{thm:transition_rate},
Chapter~\ref{chap:enzymes_topology}). Maximising the ATP yield
objective $\mathbf{c}^T\mathbf{v}$ is equivalent to maximising
$\sum_j c_j \tau_p^{-1} b^{-\Delta\Depth_j}$, which is maximised by
minimising $\sum_j c_j \Delta\Depth_j$. When $c_j$ is the stoichiometric
yield coefficient, this reduces to minimising total partition depth across
the network.
\end{proof}

\begin{remark}
This equivalence implies that the cell is performing partition depth
minimisation---a computation---when it selects among the exponentially
many feasible flux distributions. The metabolic computer is not just
a metaphor: it is computing the minimum-depth path through a
high-dimensional partition network.
\end{remark}

%------------------------------------------------------------------------------
\section{The Quantum Foundation of the Metabolic Computer}
\label{sec:quantum_metabolism}
%------------------------------------------------------------------------------

The metabolic computer is quantum at its most fundamental level. Every
partition operation has a minimum time $\tau_p = \hbar / (\kB T)$
(Chapter~\ref{chap:bounded_phase_space}). This is a quantum-mechanical
quantity: the time required for a quantum system to evolve from one
distinguishable state to an orthogonal state (Mandelstam--Tamm bound).

At $T = 310$\,K (biological temperature):
\begin{equation}
  \tau_p = \frac{\hbar}{\kB T}
         = \frac{1.055 \times 10^{-34}}{1.381 \times 10^{-23} \times 310}
         = 2.46 \times 10^{-14}\ \mathrm{s}.
  \label{eq:tau_p_bio}
\end{equation}

This sets the \emph{fundamental clock rate} of the metabolic computer:
no partition operation, no enzymatic step, can be faster than $\tau_p$.
The observed $\kcat \approx 10^6\ \mathrm{s}^{-1}$ for carbonic anhydrase
is already nine orders of magnitude slower than $\tau_p^{-1} \approx
4 \times 10^{13}\ \mathrm{s}^{-1}$, reflecting the multiple partition
boundaries that must be crossed in each catalytic cycle.

\begin{remark}
``Quantum metabolism'' does not require quantum coherence, entanglement, or
superposition of metabolite states. It requires only that the minimum
partition-transition time is set by quantum mechanics. This is a much
weaker and more defensible claim, but it has the same consequence: the
metabolic computer has a quantum-limited clock speed, and all biological
timescales are ultimately bounded below by $\tau_p$.
\end{remark}

%------------------------------------------------------------------------------
\section{Flux Tables: Healthy vs.\ Disease States}
\label{sec:flux_tables}
%------------------------------------------------------------------------------

Table~\ref{tab:metabolic_flux} gives characteristic metabolic fluxes in
healthy, insulin-resistant, and Warburg-phenotype cells.

\begin{table}[htbp]
\centering
\caption{Normalised metabolic fluxes across the five levels in three
  metabolic states. Fluxes are relative to glucose uptake $= 1.0$.
  Partition depth $\Depth_i$ is the estimated depth of the inter-level
  coupling pathway. Values derived from $^{13}$C-metabolic flux analysis
  data compiled in \citet{TeSlaa2016}.}
\label{tab:metabolic_flux}
\small
\begin{tabular}{lccccc}
\toprule
Pathway / Level & \multicolumn{2}{c}{Healthy} & \multicolumn{2}{c}{Insulin Resistant} & Warburg \\
\cmidrule(lr){2-3} \cmidrule(lr){4-5} \cmidrule(lr){6-6}
 & Flux & $\Depth_i$ & Flux & $\Depth_i$ & Flux \\
\midrule
L1: Glucose uptake          & 1.00 & 1.0 & 0.35 & 3.1 & 1.80 \\
L2: Glycolytic flux          & 1.00 & 1.2 & 0.35 & 2.8 & 1.80 \\
L2: Lactate output           & 0.05 & --- & 0.08 & --- & 1.60 \\
L3: TCA cycle flux (acetyl-CoA) & 0.95 & 1.5 & 0.27 & 2.2 & 0.20 \\
L3: TCA cycle (NADH output)  & 2.85 & --- & 0.81 & --- & 0.60 \\
L4: OxPhos flux              & 0.90 & 1.8 & 0.25 & 2.5 & 0.02 \\
L4: ATP yield (per glucose)  & 32   & --- & 22   & --- & 4    \\
L5: HIF-1$\alpha$ activity   & 0.10 & 2.1 & 0.25 & 2.8 & 0.95 \\
$D$ (hierarchical depth)     & 5.0  & --- & 3.5  & --- & 2.0  \\
$I_{\mathrm{total}}$ (bits)  & 7.29 & --- & 5.1  & --- & 2.9  \\
\bottomrule
\end{tabular}
\end{table}

The most striking feature is the ATP yield: healthy cells produce $\sim 32$
ATP per glucose, insulin-resistant cells $\sim 22$, and Warburg cells $\sim 4$.
The partition interpretation is direct: each missing level of the hierarchy
removes a layer of energy extraction. The Warburg cell has abandoned the
two most efficient levels (L3 and L4), accepting a $\sim 8$-fold reduction
in ATP yield in exchange for faster substrate processing at L1--L2.

%------------------------------------------------------------------------------
\section{Summary}
\label{sec:metabolic_summary}
%------------------------------------------------------------------------------

\begin{chapterbox}
Metabolism is a five-level hierarchical categorical computer (L1--L5)
performing information compression from the molecular to the gene-expression
timescale. The total information compressed in healthy human metabolism is
$I_{\mathrm{total}} = 7.29$ bits (Theorem~\ref{thm:compression}), measurable
from metabolomics data.

The TCA cycle (L3) is a partition loop (Theorem~\ref{thm:partition_loop}):
a cyclic partition topology satisfying $\sum_{i=1}^8 \Delta\Depth_i = 0$,
which makes oxaloacetate partition-catalytic and the cycle self-sustaining.
This is why the TCA cycle is necessarily cyclic rather than linear.

The ATP ratchet (L4) is a partition depth deficit: ATP has
$\Delta\Depth_{\mathrm{ATP}} \approx 19$--20 natural units above ADP + P$_i$,
and ATP synthase couples the proton gradient partition depth to the synthesis
of this high-$\Depth$ state.

ATP yield mirrors information hierarchy: L2 yields 2\,ATP and high information
compression (preprocessing); L3--L4 yield $\sim 28$--$30$\,ATP and the deepest
partition operations (Proposition~\ref{prop:atp_information}).

Hierarchical depth $D$ is a health metric: $D = 5$ (healthy, 7.29\,bits),
$D \approx 4.8$ (lithium-stabilised bipolar, 7.0\,bits), $D \approx 4.2$
(metformin-treated T2D, 6.1\,bits), $D \approx 3$--$3.5$ (insulin resistance),
$D \approx 2$ (Warburg effect, 2.9\,bits). Cascade failure propagates upward:
L3 failure forces L4 and L5 failure (Theorem~\ref{thm:cascade_failure}).

Pharmaceutical interventions (metformin at L4, lithium at L5, DCA at L3--L4)
are topologically targeted partition depth corrections quantified by the
efficacy formula $\eta = \Delta D / (5 - D_{\mathrm{pre}})$.

The equivalence of FBA and partition depth minimisation
(Theorem~\ref{thm:fba_equivalence}) establishes that the metabolic computer
is performing genuine optimisation, bounded below in clock speed by the
quantum partition time $\tau_p = \hbar/(\kB T) = 2.46 \times 10^{-14}$\,s
at biological temperature.
\end{chapterbox}
